% Original academic essay starter based on the report-oriented use case of
% Cinofix's LaTeX scientific template collection.
\documentclass[11pt,a4paper]{article}
\usepackage[margin=2.7cm]{geometry}
\usepackage{amsmath}
\usepackage{booktabs}
\usepackage{microtype}
\usepackage{hyperref}

\title{Reproducible Evaluation of a Small Data Pipeline}
\author{Your Name\\Department or Institution}
\date{\today}

\begin{document}
\maketitle

\begin{abstract}
This essay-style starter shows a compact scientific workflow: define a claim,
describe a measurable protocol, report results in a table, and state the
limits of the evidence. Replace the sample study with a project of your own.
\end{abstract}

\section{Question and claim}
Engineering and scientific projects often begin with a plausible improvement
that has not yet been measured under repeatable conditions. This paper asks
whether batching independent records reduces the time needed to compute a
summary while preserving the result. The claim is deliberately narrow: a
batching strategy is useful only when its speedup is larger than the cost of
preparing and validating the batches.

\section{Method}
The evaluation compares a sequential baseline with a batched implementation.
Both variants process the same fixed input and write a checksum after each
run. If $T_0$ is the baseline time and $T_1$ the batched time, the relative
speedup is
\begin{equation}
  S=\frac{T_0}{T_1}.
\end{equation}
The experiment is repeated five times after one warm-up run. The machine,
input size, software version, and random seed are recorded with each result.

\section{Results}
Table~\ref{tab:results} shows a compact example of the evidence record. The
checksums agree in every run, so the comparison concerns execution time rather
than a changed output.

\begin{table}[ht]
  \centering
  \caption{Illustrative repeated measurements.}
  \label{tab:results}
  \begin{tabular}{lrrr}
    \toprule
    Run & Baseline (ms) & Batched (ms) & Speedup \\
    \midrule
    1 & 842 & 516 & 1.63 \\
    2 & 835 & 509 & 1.64 \\
    3 & 849 & 514 & 1.65 \\
    4 & 838 & 511 & 1.64 \\
    5 & 844 & 512 & 1.65 \\
    \bottomrule
  \end{tabular}
\end{table}

\section{Discussion}
The example results support the limited claim: batching reduces the measured
runtime for this input while preserving the checksum. They do not show that
the strategy is always faster. Small inputs may be dominated by setup cost,
and a different distribution of records may make memory pressure the limiting
factor. A useful scientific template makes those conditions visible alongside
the headline number.

\section{Conclusion}
The evaluation should be extended with larger inputs, a cold-cache run, and a
failure-injection test. Until those tests are complete, the responsible
conclusion is that batching is promising for the measured workload, not that
it is a universal optimization.

\end{document}
